% !TEX root = master.tex
\chapter{Referenzdokument}
\label{ch:a-referenz}

\section{Ziele}
\emph{GuardX} ist ein modulares \enquote{Web Security \& Privacy Toolkit}. Ziel des Projekts ist
eine Web-App, die registrierten Nutzerinnen und Nutzern mehrere praxisnahe Werkzeuge rund um
digitale Sicherheit und Privatsph\"are an einer Stelle bereitstellt und dabei selbst nach
g\"angigen Sicherheitsprinzipien aufgebaut ist. Die App verfolgt damit einen doppelten Zweck:
Zum einen liefert sie konkreten Nutzen (Passw\"orter pr\"ufen, Bild-Metadaten entfernen,
Tracking-Risiken verstehen), zum anderen demonstriert sie die im Modul geforderten technischen
Konzepte -- insbesondere Sessionverwaltung, sichere Authentifizierung, Rechteverwaltung und
Datenbankzugriffe.

\section{Funktionsrahmen}
Die Anwendung gliedert sich in einen \textbf{Account- und Verwaltungsbereich} und drei
\textbf{Fachmodule}.

Der Account- und Verwaltungsbereich umfasst:
\begin{itemize}
  \item Registrierung, Login und Logout mit serverseitiger Validierung,
  \item Sessionverwaltung mit automatischem Logout nach konfigurierbarer Inaktivit\"at,
  \item einen Admin-Bereich mit vollst\"andiger Benutzerverwaltung (\ac{crud}),
  \item Passwortverwaltung (\"Andern durch Nutzer, Zur\"ucksetzen durch Admin) und
  \item eine Rechteverwaltung mit den Rollen \emph{Admin} und \emph{User}.
\end{itemize}

Die drei Fachmodule sind:
\begin{itemize}
  \item \textbf{Metadaten-Bereiniger}: Entfernt \ac{exif}-Daten aus hochgeladenen Bildern und
        bietet die ausgelesenen Metadaten als \ac{json}-Datei zum Download an.
  \item \textbf{Browser-Fingerprinting}: Macht sichtbar, welche Merkmale (Canvas, WebGL,
        Zeitzone, Sprache, Hardware) einen Browser identifizierbar machen.
  \item \textbf{Live-Passwort-Checker}: Pr\"uft per \ac{ajax} in Echtzeit die Passwortst\"arke und
        gleicht das Passwort datenschutzfreundlich (k-Anonymity) gegen bekannte Datenlecks ab.
\end{itemize}

Querschnittlich werden die geforderten Technologien eingesetzt: client- und serverseitige
Eingabepr\"ufung (u.\,a. mit regul\"aren Ausdr\"ucken), Einbindung von jQuery/jQuery~UI und
Bootstrap, dynamisches Nachladen via \ac{ajax}, Meldungs- und Best\"atigungsdialoge, Datenexport
via \ac{json} sowie eine dynamische, responsive Men\"ustruktur. Eine vollst\"andige Zuordnung
dieser Funktionen zu den jeweiligen Skripten findet sich in Kapitel~\ref{ch:a-funktionen}.

\section{Projektrahmen}
Das Projekt wurde als zeitlich begrenzte Teamarbeit von vier Personen im Sommersemester~2026
umgesetzt. Die technologischen Rahmenbedingungen entsprechen den Projektvorgaben:

\begin{table}[htbp]
\centering
\caption{Technologischer Rahmen von Projekt~A}
\label{tab:a-rahmen}
\begin{tabularx}{\textwidth}{@{}lX@{}}
\toprule
\textbf{Bereich} & \textbf{Technologie} \\
\midrule
Client       & HTML5, \ac{css}, JavaScript; jQuery, jQuery~UI, Bootstrap~5.3 \\
Server       & Apache, \ac{php}~8.2.4 (ohne Framework) \\
Datenbank    & MariaDB~10.4.28 \\
Laufzeitumgebung & XAMPP \\
Versionsverwaltung & Git / GitHub \\
\bottomrule
\end{tabularx}
\end{table}

Auf den Einsatz fertiger Frameworks (z.\,B. Angular, Vue, Symfony, React) wurde vorgabengem\"a\ss{}
verzichtet; eingesetzt werden ausschlie\ss{}lich die erlaubten Bibliotheken jQuery, jQuery~UI und
Bootstrap.

\section{Umsetzungsstrategie}
Die Entwicklung erfolgte iterativ und arbeitsteilig. Zur Zusammenarbeit wurde ein
Git-Feature-Branch-Workflow genutzt: Neue Funktionen wurden in eigenen Branches entwickelt,
per Pull-Request zusammengef\"uhrt und vor dem Merge daraufhin gepr\"uft, ob die Anwendung unter
XAMPP lauff\"ahig bleibt. Wiederkehrende Bestandteile (Header/Navigation, Footer, Konfiguration)
wurden als wiederverwendbare Includes ausgelagert, sodass alle Seiten ein einheitliches
Erscheinungsbild und Verhalten besitzen. Sicherheitsrelevante Mechanismen
(\ac{csrf}-Schutz, Passwort-Hashing, Prepared Statements) wurden von Beginn an mitgedacht und
sind zentral umgesetzt. Vor der Abgabe wurde die Installierbarkeit auf einem frischen System
getestet.

\section{Aufgaben und Zust\"andigkeiten}
Die Arbeitspakete wurden wie in Tabelle~\ref{tab:a-zustaendig} dargestellt verteilt. Die
Zuordnung orientiert sich an den jeweils verantworteten Komponenten; \"ubergreifende Aufgaben
(Code-Reviews, Integration, Test) wurden gemeinsam wahrgenommen.

\begin{table}[htbp]
\centering
\caption{Aufgaben und Zust\"andigkeiten in Projekt~A}
\label{tab:a-zustaendig}
\begin{tabularx}{\textwidth}{@{}lX@{}}
\toprule
\textbf{Teammitglied} & \textbf{Verantwortungsbereich} \\
\midrule
Maik Baumg\"artner & Backend: Session- und Benutzerverwaltung, Login/Logout/Registrierung,
  Admin-Bereich, Datenbankanbindung und -schema, zentrale Hilfsfunktionen sowie
  Navigations-Templates. \\
\addlinespace
Marcel Reinholz & Footer und alle dar\"uber verlinkten Rechtsseiten (Impressum, Datenschutz,
  AGB) sowie das Modul der API-Aufrufe (Live-Passwort-Checker, Leak-Abgleich). \\
\addlinespace
Jakob Goelden & Modul Browser-Fingerprinting (\texttt{info.php}) sowie Einrichtung und
  Pflege des GitHub-Repositories (Branch-/Merge-Strategie, Versionsverwaltung). \\
\addlinespace
Moritz Brenner & Modul Metadaten-Bereiniger, das durchg\"angige \ac{css}-Design
  (Layout, Cyber-Theme, responsive Anpassungen) sowie die Startseite. \\
\bottomrule
\end{tabularx}
\end{table}
